\section*{化学}

\subsection*{卤素单质性质}
F$_2$：\sbl 色 \sbl 体 \quad 
Cl$_2$：\sbl 色 \sbl 体

Br$_2$：\sbl 色 \sbl 体 \quad 
I$_2$：\sbl 色 \sbl 体

\subsection*{化合物特点}
离子化合物特点：\mbl \quad 
共价化合物特点：\mbl \quad 
极性化合物：\mbl

离子、共价化合物的电子式、结构式、化学式，各举一例：
\nl{3}

球棍模型、VSEPR 模型、电子云模型、比例模型，各举一例：
\nl{3}

\subsection*{原子结构与周期表}
如何区分 s/p/d/f 层先后顺序：

元素周期表族名称顺序：

默写周期表21-30位（名称与符号）：
\nl{2}
\subsection*{VSEPR 理论}
VSEPR模型：中心原子连

\column{
两个原子：\mbl \sep
三个原子：\mbl \sep
四个原子：\mbl
}

\column{
两个原子和一个孤电子对：\mbl \sep
两个原子和两个孤电子对：\mbl
}

三个原子和一个孤电子对：\mbl

分子几何构型：中心原子连

\column{
两个原子：\mbl \sep
三个原子：\mbl \sep
四个原子：\mbl
}

\column{
两个原子和一个孤电子对：\mbl \sep
两个原子和两个孤电子对：\mbl
}

三个原子和一个孤电子对：\mbl

\subsection*{配位化学}
配位键是：

配位数是：

晶体具有各向 \sbl 性，熔化后呈现 \sbl 型，非晶体则为 \sbl 型

常见配合物：

检验铁离子用 \sbl ，成 \sbl 色

检测醛基或溶解难溶性银盐用 \sbl ， \sbl 于水

铜离子加氨水，先生成 \sbl ，过量生成 \sbl 融性 \sbl 色 \sbl

\subsection*{化学平衡}
化学平衡常数（气体、液体）：\lbl

沉淀溶解平衡计算：\lbl

字母含义：

\column{
$K_a$\sbl \sep
$K_b$\sbl \sep
$K_w$\sbl \sep
$K_p$\sbl
}
\column{
$K_h$\sbl \sep
$K_{sp}$\sbl \sep
$K_c$\sbl
}

\subsection*{官能团}
\column{
碳卤键/ \sbl \sep
氰基/ \sbl \sep
醚键/ \sbl
}
\nl{2}
\column{
醛基/ \sbl \sep
酮羰基/ \sbl \sep
酯基/ \sbl
}
\nl{2}
\column{
氨基/ \sbl \sep
羧基/ \sbl \sep
酰胺基/ \sbl
}
\nl{2}
\subsection*{命名与谱学}
基团命名顺序：

\sbl $>$ \sbl $>$ \sbl $>$ \sbl $>$ \sbl $>$ \sbl $>$ \sbl $>$ \sbl \\
\sbl $>$ \sbl $>$ \sbl

\column{
质谱法测：\sbl \sep
红外光谱测：\sbl \sep
核磁氢谱测：\sbl \sep
X 射线衍射测：\sbl
}

\subsection{常见元素及其化合物}

\subsubsection{颜色与还原性}

\column{\ce{}\csb \sep \ce{NO2}\csb \sep \ce{Br2}\csb \sep \ce{}\csb \sep \ce{}\csb }
\column{\ce{}\csb \sep \ce{}\csb \sep \ce{ Cr2O7^{2-}}\csb \sep \ce{Cr^{2+}}\csb \sep \ce{}\csb }
% \column{\ce{ }\csb \sep \ce{ }\csb \sep \ce{ }\csb \sep \ce{ }\csb \sep \ce{ }\csb }

还原性：

$\sbl > \sbl > \sbl > \sbl > \sbl > \sbl > \sbl$

氧化性：

$\sbl > \sbl > \sbl > \sbl > \sbl > \sbl > \sbl$

\subsubsection{Na}
\column{颜色\csb \sep 软硬\csb \sep 密度 \csb}

与以下是否反应：

\column{O$_2$ \sep X$_2$}
\nl{1}
\column{H$_2$O \sep CuSO$_4$}
\nl{1}
\column{C$_2$H$_5$OH \sep}
\nl{1}
\column{NaK合金性质\cmb \sep 高压钠灯颜色\csb 穿透性\csb}

Na$_2$O与Na$_2$O$_2$颜色、能和什么反应，生成物区别：

NaHCO$_3$“小”在哪里：

侯氏制碱法流程：
\nl{1}

\ce{NaNO2}可被用作\mbl， 毒性\sbl 。与以下是否反应：

\column{NH$_4$ \sep KMnO$_4$}

\ce{NaH}是\sbl 化合物，其各元素化合价为\mbl

\subsection{Fe}

与以下是否反应：

\column{\ce{Cl2} \sep \ce{S}}
\nl{1}
\column{\ce{H2O} \sep \ce{HNO3}}
\nl{1}
\column{\ce{Cl2} \sep \ce{S}}
\nl{1}
\column{\ce{Fe2O3 + HI} \sep \ce{FeO + HNO3}}
\nl{1}

\ce{Fe2+}：

与以下是否反应，有什么现象：

\column{\ce{MnO4-} \sep \ce{Cr2O7^{2-}}}
\nl{1}
\column{\ce{HNO3} \sep \ce{H2O2}}
\nl{1}
\column{\ce{H2O} \sep \ce{CO3^{2-}}}
\nl{1}
\column{\ce{HCO3-} \sep \ce{KSCN}}
\nl{1}
\column{\ce{K3[Fe(CN)6]} \sep \ce{NaOH}}
\nl{1}

\ce{Fe}或\ce{FeO}变成\ce{Fe3O4}：
\nl{1}

\ce{Fe^{3+}}：
水中\sbl 色，可还原\sbl / \sbl / \sbl / \sbl / \sbl ，\sbl 可腐蚀铜电路板

与以下是否反应：

\ce{H2O}，完全沉淀pH值\csb 配置时添加\sbl

\nl{1}
\column{\ce{SCN-} \sep \ce{NaOH}}
\nl{1}

\ce{K2FeO4}：氧化性\sbl 可用于\mbl、\mbl、\mbl

